% Préambule des annexes : macros, styles TikZ, couleurs et références repris tels quels de la soutenance
% de thèse (ludwig-h/manuscrit-de-th-se, Soutenance/soutenance/main.tex).
\colorlet{inria-bleu}{inria-2024-bleu-azur}
% ------------------------------------------------------------
% Système de citations entre crochets.
%   \DeclareRef{cle}{Label}{Référence complète}   — littérature
%   \DeclareMyRef{cle}{Label}{Référence complète} — mes travaux (rouge)
%   \citb{cle}    : [Label] dans le corps du slide
%   \reffoot{c1,c2} : références complètes en bas du slide
%                     (à placer sur le slide de première citation)
% ------------------------------------------------------------
\makeatletter
\newcommand{\DeclareRef}[3]{%
  \expandafter\gdef\csname ref@lab@#1\endcsname{#2}%
  \expandafter\gdef\csname ref@full@#1\endcsname{#3}}
\newcommand{\DeclareMyRef}[3]{\DeclareRef{#1}{#2}{#3}%
  \expandafter\gdef\csname ref@mine@#1\endcsname{}}
\newcommand{\refbracket}[1]{%
  \ifcsname ref@mine@#1\endcsname
    \textcolor{inria-rouge}{[\csname ref@lab@#1\endcsname]}%
  \else
    \textcolor{gris_fonce_inria}{[\csname ref@lab@#1\endcsname]}%
  \fi}
\newcommand{\citb}[1]{\mbox{{\footnotesize\bfseries\refbracket{#1}}}}
\newcommand{\bibligne}[1]{\item {\bfseries\refbracket{#1}}\ \csname ref@full@#1\endcsname}
% Le \par de chaque référence doit tomber DANS le groupe : sinon la dernière
% ligne est coupée avec le \baselineskip du corps, et non celui de \tiny —
% ce qui intercalait une demi-ligne vide avant la dernière référence.
\newcommand{\reffoot}[1]{%
  \begin{textblock}{0.90}[0,1](0.05,0.925)%
    \begingroup
    \tiny\linespread{0.90}\selectfont\raggedright\color{gris_fonce_inria}%
    \setlength{\parindent}{0pt}\setlength{\parskip}{0pt}%
    \edef\rf@list{\zap@space#1 \@empty}% tolère « a, b » comme « a,b »
    \@for\rf@key:=\rf@list\do{%
      \noindent{\bfseries\expandafter\refbracket\expandafter{\rf@key}}%
      \space\csname ref@full@\rf@key\endcsname\par}%
    \endgroup
  \end{textblock}}
\makeatother

% ---- Littérature -------------------------------------------
\DeclareRef{chazal13}{ToMATo 13}{F.~Chazal, L.~J.~Guibas, S.~Y.~Oudot, P.~Skraba, « Persistence-based clustering in Riemannian manifolds » (ToMATo), \textit{J. ACM}, 60(6), art.~41, 2013.}
\DeclareRef{ghrist08}{Ghrist 08}{R.~Ghrist, « Barcodes: the persistent topology of data », Bulletin of the AMS, 2008.}
\DeclareRef{hartigan81}{Hartigan 81}{J.~A.~Hartigan, « Consistency of Single Linkage for high-density clusters », JASA, 1981.}
\DeclareRef{biau15}{Biau--Devroye 15}{G.~Biau, L.~Devroye, \textit{Lectures on the Nearest Neighbor Method}, Springer, 2015.}
\DeclareRef{florek51}{Florek 51}{K.~Florek, J.~Łukaszewicz, J.~Perkal, H.~Steinhaus, S.~Zubrzycki, « Sur la liaison et la division des points d'un ensemble fini », Colloq. Math., 1951.}
\DeclareRef{gowerross69}{Gower--Ross 69}{J.~C.~Gower, G.~J.~S.~Ross, « Minimum spanning trees and Single-Linkage cluster analysis », JRSS C, 1969.}
\DeclareRef{penrose03}{Penrose 03}{M.~Penrose, \textit{Random Geometric Graphs}, Oxford University Press, 2003.}
\DeclareRef{meesterroy96}{Meester--Roy 96}{R.~Meester, R.~Roy, \textit{Continuum Percolation}, Cambridge Tracts in Mathematics 119, Cambridge University Press, 1996.}
\DeclareRef{cameron44}{Cameron--Martin 44}{R.~H.~Cameron, W.~T.~Martin, « Transformations of Wiener integrals under translations », \textit{Annals of Mathematics}, 45(2), 1944.}
\DeclareRef{schilder66}{Schilder 66}{M.~Schilder, « Some asymptotic formulas for Wiener integrals », \textit{Transactions of the AMS}, 125(1), 1966.}
\DeclareRef{dembozeitouni10}{Dembo--Zeitouni 10}{A.~Dembo, O.~Zeitouni, \textit{Large Deviations Techniques and Applications}, 2\up{e} éd., Springer, 2010.}
\DeclareRef{berlinet04}{Berlinet--Thomas-Agnan 04}{A.~Berlinet, C.~Thomas-Agnan, \textit{Reproducing Kernel Hilbert Spaces in Probability and Statistics}, Kluwer, 2004.}
\DeclareRef{boissonnat18}{Boissonnat 18}{J.-D.~Boissonnat, F.~Chazal, M.~Yvinec, \textit{Geometric and Topological Inference}, Cambridge University Press, 2018.}
\DeclareRef{chaudhuri10}{Chaudhuri--Dasgupta 10}{K.~Chaudhuri, S.~Dasgupta, « Rates of convergence for the cluster tree », \textit{Advances in Neural Information Processing Systems} (NIPS), vol.~23, p.~343--351, 2010.}
\DeclareRef{ester96}{Ester 96}{M.~Ester, H.-P.~Kriegel, J.~Sander, X.~Xu, « A density-based algorithm for discovering clusters » (DBSCAN), KDD, 1996.}
\DeclareRef{campello13}{Campello 13}{R.~Campello, D.~Moulavi, J.~Sander, « Density-based clustering based on hierarchical density estimates » (HDBSCAN), PAKDD, 2013.}
\DeclareRef{callahan95}{Callahan--Kosaraju 95}{P.~B.~Callahan, S.~R.~Kosaraju, « A decomposition of multidimensional point sets with applications to $k$-nearest-neighbors and $n$-body potential fields », \textit{J. ACM}, 42(1), 1995.}
\DeclareRef{edelsbrunner23}{Edelsbrunner--Osang 23}{H.~Edelsbrunner, G.~Osang, « A simple algorithm for higher-order Delaunay mosaics and alpha shapes », \textit{Algorithmica}, 2023.}
\DeclareRef{forina83}{Forina 83}{M.~Forina, C.~Armanino, S.~Lanteri, E.~Tiscornia, « Classification of olive oils from their fatty acid composition », 1983 (jeu de données).}
\DeclareRef{rolle23}{Rolle--Scoccola 24}{A.~Rolle, L.~Scoccola, « Stable and consistent density-based clustering via multiparameter persistence » (Persistable), \textit{JMLR}, 25(258), 2024.}
\DeclareRef{semantickitti19}{SemanticKITTI 19}{J.~Behley \emph{et al.}, « SemanticKITTI: A Dataset for Semantic Scene Understanding of LiDAR Sequences », \textit{ICCV}, 2019.}
\DeclareRef{sautier25}{ALPINE 25}{C.~Sautier, G.~Puy, A.~Boulch, R.~Marlet, V.~Lepetit, « Is clustering enough for LiDAR instance segmentation? A state-of-the-art training-free baseline » (ALPINE), 2025.}
\DeclareRef{oh25}{Geo-4D 25}{G.~Oh, Y.~Jang, J.~Choi, S.-J.~Kang, G.~Lin, S.~Kim, « Training-free global geometric association for 4D LiDAR panoptic segmentation » (Geo-4D), 2025.}
\DeclareRef{avrachenkov22}{Avrachenkov--Dreveton 22}{K.~Avrachenkov, M.~Dreveton, \textit{Statistical Analysis of Networks}, Now Publishers, 2022.}
\DeclareRef{sw87}{Swendsen--Wang 87}{R.~H.~Swendsen, J.-S.~Wang, « Nonuniversal critical dynamics in Monte Carlo simulations », \textit{Physical Review Letters}, 1987.}
\DeclareRef{es88}{Edwards--Sokal 88}{R.~G.~Edwards, A.~D.~Sokal, « Generalization of the Fortuin--Kasteleyn--Swendsen--Wang representation and Monte Carlo algorithm », \textit{Physical Review D}, 38(6), 1988.}
\DeclareRef{sb18}{Sankararaman--Baccelli 18}{A.~Sankararaman, F.~Baccelli, « Community detection on Euclidean random graphs », \textit{ACM--SIAM SODA}, 2018.}
\DeclareRef{frangi98}{Frangi 98}{A.~F.~Frangi, W.~J.~Niessen, K.~L.~Vincken, M.~A.~Viergever, « Multiscale vessel enhancement filtering », MICCAI, 1998.}
\DeclareRef{kleinberg02}{Kleinberg 02}{J.~Kleinberg, « An impossibility theorem for clustering », \textit{Advances in Neural Information Processing Systems} (NIPS), vol.~15, MIT Press, p.~446--453, 2002.}
\DeclareRef{carlsson10}{Carlsson--Mémoli 10}{G.~Carlsson, F.~Mémoli, « Characterization, stability and convergence of hierarchical clustering methods », JMLR, 2010.}
\DeclareRef{culbertson18}{Culbertson 18}{J.~Culbertson, D.~P.~Guralnik, P.~F.~Stiller, « Functorial hierarchical clustering with overlaps », \textit{Discrete Applied Mathematics}, 2018.}
\DeclareRef{brown20}{GPT-3 20}{T.~Brown \emph{et~al.}, « Language models are few-shot learners », \textit{NeurIPS}, 2020.}
\DeclareRef{robert23}{SPT 23}{D.~Robert, H.~Raguet, L.~Landrieu, « Efficient 3D semantic segmentation with Superpoint Transformer », \textit{ICCV}, 2023.}
\DeclareRef{ravi24}{SAM 2 24}{N.~Ravi \emph{et~al.}, « SAM 2: segment anything in images and videos », 2024.}
\DeclareRef{sonata25}{Sonata 25}{X.~Wu \emph{et~al.}, « Sonata: self-supervised learning of reliable point representations », \textit{CVPR}, 2025.}
\DeclareRef{concerto25}{Concerto 25}{Y.~Zhang \emph{et~al.}, « Concerto: joint 2D--3D self-supervised learning emerges spatial representations », \textit{NeurIPS}, 2025.}
\DeclareRef{yu25}{PolyhedronNet 25}{D.~Yu, G.~Zhang, L.~Zhao, « PolyhedronNet: representation learning for polyhedra with surface-attributed graph », \textit{ICLR}, 2025.}
\DeclareRef{hsa25}{HSA 25}{S.~Amizadeh \emph{et~al.}, « Hierarchical self-attention: generalizing neural attention mechanics to multi-scale problems », \textit{NeurIPS}, 2025.}
\DeclareRef{utonia26}{Utonia 26}{Y.~Zhang \emph{et~al.}, « Utonia: toward one encoder for all point clouds », 2026.}
\DeclareRef{ayana25}{AYANA RA 25}{Équipe-projet \textsc{Ayana}, « 3D LiDAR Instance Segmentation with Geometric Priors: The HGP-Clusterer », \href{https://radar.inria.fr/report/2025/ayana/index.html\#AYANA-RA-2025_label_NavalResult}{\textit{Rapport annuel d'activité Inria 2025}}, \S~7.7.}
\DeclareRef{bommasani21}{Bommasani 21}{R.~Bommasani \emph{et~al.}, « On the opportunities and risks of foundation models », 2021.}
\DeclareRef{kirillov23}{SAM 23}{A.~Kirillov \emph{et~al.}, « Segment Anything », \textit{ICCV}, 2023.}
\DeclareRef{hu22}{LoRA 22}{E.~J.~Hu \emph{et~al.}, « LoRA: low-rank adaptation of large language models », \textit{ICLR}, 2022.}
\DeclareRef{ge24}{CrackSAM 24}{K.~Ge \emph{et~al.}, « Fine-tuning vision foundation model for crack segmentation in civil infrastructures », \textit{Construction and Building Materials}, 2024.}
\DeclareRef{ying21}{Graphormer 21}{C.~Ying \emph{et~al.}, « Do Transformers really perform badly for graph representation? », \textit{NeurIPS}, 2021.}
\DeclareRef{turaga09}{MALIS 09}{S.~C.~Turaga, K.~L.~Briggman, M.~Helmstaedter, W.~Denk, H.~S.~Seung, « Maximin affinity learning of image segmentation », \textit{NIPS}, 2009.}

% ---- Mes travaux (en rouge) --------------------------------
\DeclareMyRef{hal17}{HAL 17}{L.~Hauseux (encadré par B.~Błaszczyszyn), « Un classificateur non supervisé utilisant les complexes simpliciaux », Inria Paris, 2017. HAL : hal-01597846.}
\DeclareMyRef{cn23}{CN 23}{L.~Hauseux, K.~Avrachenkov, J.~Zerubia, « Graph based approach for galaxy filament extraction », \textit{Complex Networks \& Their Applications}, Menton, 2023.}
\DeclareMyRef{eusipco24}{EUSIPCO 24}{L.~Hauseux, K.~Avrachenkov, J.~Zerubia, « Benefits of hypergraphs for density-based clustering », \textit{EUSIPCO}, Lyon, 2024.}
\DeclareMyRef{cn24}{CN 24}{L.~Hauseux, K.~Avrachenkov, J.~Zerubia, « Hypergraphs, percolation and hierarchical clustering », \textit{Complex Networks \& Their Applications}, Istanbul, 2024.}
\DeclareMyRef{src24}{SIGMETRICS SRC 24}{L.~Hauseux, « How can we theoretically measure the performance of density-based clustering algorithms? », ACM SIGMETRICS SRC, Venise, 2024 (2\up{e} prix).}
\DeclareMyRef{ans26}{ANS 26}{L.~Hauseux, K.~Avrachenkov, J.~Zerubia, « Generalization of Single-Linkage with higher-order interactions », \textit{Applied Network Science}, 11(9), 2026.}
\DeclareMyRef{gretsi25}{GRETSI 25}{L.~Hauseux, R.~Antoine, P.~Foucher, P.~Charbonnier, J.~Zerubia, « Généralisation du filtre de Frangi pour l'extraction des réseaux de fissures », GRETSI, 2025.}
\DeclareMyRef{asilomar25}{Asilomar 25}{L.~Hauseux, N.~Soprano-Loto, K.~Avrachenkov, « Higher-order Monte Carlo cluster dynamics for community detection in Euclidean graphs », Asilomar, 2025.}
\DeclareMyRef{ii26}{I\&I 26+}{L.~Hauseux, N.~Soprano-Loto, K.~Avrachenkov, « A Bayesian framework for community detection on signed graphs », article en cours de rédaction pour \textit{Information and Inference}, 2026.}
\DeclareMyRef{euvip26}{EUVIP 26}{L.~Hauseux, R.~Antoine, P.~Foucher, P.~Charbonnier, J.~Zerubia, « Multi-modal, training-free crack extraction via generalized Frangi graph », EUVIP, 2026.}
\DeclareMyRef{isprs26}{ISPRS 26+}{L.~Hauseux, R.~Antoine, P.~Foucher, P.~Charbonnier, J.~Zerubia, article en préparation pour l'\textit{ISPRS Journal of Photogrammetry and Remote Sensing}, 2026.}

% ------------------------------------------------------------
% Styles TikZ (repris de la présentation NIM / manuscrit)
% ------------------------------------------------------------
\tikzset{
  chapbox/.style={draw,rounded corners,very thick,align=center,fill=gray!8,inner sep=5pt,minimum height=1.05cm},
  chapboxdark/.style={draw,rounded corners,very thick,align=center,fill=gray!18,inner sep=5pt,minimum height=1.05cm},
  chaparr/.style={-{Latex[length=2.4mm]},very thick},
  chapdash/.style={-{Latex[length=2.4mm]},very thick,dashed},
  spin/.style={circle,draw,inner sep=0pt,minimum size=5mm,font=\scriptsize,fill=white},
  tinynode/.style={circle,draw,inner sep=0pt,minimum size=3.8mm,font=\tiny,fill=white},
  posedge/.style={line width=1pt},
  negedge/.style={line width=1pt,dashed},
  freeze/.style={line width=1.6pt},
  faint/.style={line width=.45pt,gray!55},
  read/.style={line width=5pt, gray!45, line cap=round, preaction={draw, white, line width=7pt, line cap=round}},
  readlink/.style={densely dotted, thick, gray!60},
  errlabel/.style={font=\scriptsize\bfseries, text=red!80!black, fill=white, inner sep=1.5pt, rounded corners=2pt},
  interedge/.style={negedge, draw=gray!60}
}
\newcommand{\ind}[1]{\mathbb{I}_{\{#1\}}}


% Couleurs matplotlib (annotations des courbes de percolation)
\definecolor{mplC0}{HTML}{1F77B4}
\definecolor{mplC1}{HTML}{FF7F0E}
\definecolor{mplC2}{HTML}{2CA02C}
\definecolor{mplC3}{HTML}{D62728}
\definecolor{mplC4}{HTML}{9467BD}

% Encadré des formules importantes
\newcommand{\formulebox}[1]{%
  {\setlength{\fboxsep}{8pt}\setlength{\fboxrule}{0.8pt}%
   \fcolorbox{inria-rouge}{gray!4}{#1}}}
